\section{Scope, Limitations, and What We Do Not Claim}
\label{sec:discussion}

\textbf{Scale.} Every model in this paper is under 200K parameters
(170,896--185,000 across the tested $K$'s); this is a toy-scale
mechanistic study of one architecture's training objective, not a
scaling-law paper. No result here transfers to larger models, other
tokenizations, or the operator-count ($R$) axis without a fresh check --
the trainability fix and its replication are single-relation ($R{=}1$)
or $R{=}3$ only, never both, and no benchmark comparison against
published length/depth-generalization baselines is offered.
% evidence: NOVEL_ARCH_WATERFALL.md S9.3/S11.4a Coverage

\textbf{We do not claim the K-wall is solved.} $K{=}32$ is a closed
failure at every dimension and budget tested under the training recipe
used through Section~\ref{sec:wall} -- twenty-four independently
re-verified cells, far-depth recovery exactly $0.0000$ in every one.
The diagnosis (Section~\ref{sec:diagnosis}) explains \emph{why} a
post-hoc orthogonal projection recovers depth on an already-trained
operator; it does not, by itself, establish that training the constraint
directly will converge or generalize past $K{=}24$ -- that is exactly
what Section~\ref{sec:pending}'s now-landed run tested, and the answer
was negative: the constrained write did not train (FAIL, conditional on
the code re-audit disclosed there). We report the odds we assigned
before running it, not after.

\textbf{We do not claim a new state-tracking expressivity result.}
Section~\ref{sec:related} names two published results
\citep{grazzi2025unlocking,peng2025rwkv7} that already show matrix-valued
recurrent state can escape the transformer expressivity ceiling. NCR's
contribution is a query-time access-complexity claim, $O(\log h)$ against
every published alternative's $O(h)$ sequential rollout, with exactness
as its empirical companion -- not a claim that matrix state can do
something no architecture has been shown to do.

\textbf{We do not claim the orthogonality-constrained write is
unclaimed territory.} Section~\ref{sec:diagnosis} discloses MuonSSM
\citep{nguyen2026muonssm} as closely related, recently published, and
independently motivated prior art on orthogonalizing fast-weight content
rather than the recurrent transition matrix; our differentiation (full-rank
vs.\ rank-1, near-exact vs.\ single-step, compositional-depth-motivated vs.\
general-stability-motivated) is stated plainly, not implied by omission.

\textbf{We do not claim the mechanism is the only cause of the wall.}
Leakage-shape analysis (Section~\ref{sec:wall}) is scoped to $K \in
\{16, 24\}$ only and does not explain why $K{=}32$ fails at
\emph{every} tested dimension, including the tight-spare convention that
still works at $K \leq 24$; an absolute-$K$ factor beyond the
spare-fraction story is left open. The training-budget ANOMALY (write
residual non-monotonic across budget, observed independently at $K{=}16$
and $K{=}32$) has no loss-curve signature under the one check performed
and is reported as an open, twice-observed phenomenon, not resolved.
% evidence: NOVEL_ARCH_WATERFALL.md S11.5/S11.6

% NOTE: resolve after design-doc §10 (bug-vs-mechanism code re-audit) lands -- this passage's "pending fix landing NULL" framing awaits that resolution.
\textbf{Falsification, restated.} The bound half of this paper's claim --
that the $K{=}32$ wall is not rescued by dimension convention or training
budget alone under the recipe tested -- would be voided by any single
untested lever (a fourth budget multiple, a different anneal schedule, a
different loss) reaching CONVERGED-ROBUST with median far-depth recovery
at $h^\ast \geq 0.9$; none of the twenty-four cells tested against the
levers we did try did so. The mechanism half would be voided by the
pending fix (Section~\ref{sec:pending}) landing NULL with a clean
mechanistic signature (departure-from-normality and condition number
already near their healthy targets) -- that specific combination would
mean the diagnosis is wrong, not merely incomplete, and we would report
it as such.
